% ============================================================
% Section 5: Results
% ============================================================

% --------------------------------------------------------
\subsection{Primary Result: Patch Generation Rate}
\label{sec:results:primary}
% --------------------------------------------------------

Table~\ref{tab:main_results} reports the primary quantitative results across
all 300 SWE-bench Lite instances.

\begin{table}[t]
\centering
\caption{Primary results on SWE-bench Lite (300 instances). H: hierarchical
system; B: single-agent baseline. Discordant pairs drive the McNemar test.}
\label{tab:main_results}
\begin{tabular}{lcc}
\toprule
\textbf{Outcome} & \textbf{H} & \textbf{B} \\
\midrule
Patches generated & 53 (17.7\%) & 24 (8.0\%) \\
\midrule
H-only patches (wins) & \multicolumn{2}{c}{49} \\
B-only patches (wins) & \multicolumn{2}{c}{20} \\
Both patched          & \multicolumn{2}{c}{4}  \\
Neither patched       & \multicolumn{2}{c}{227} \\
\midrule
McNemar $\chi^2$ & \multicolumn{2}{c}{11.36} \\
$p$-value         & \multicolumn{2}{c}{$6.0 \times 10^{-4}$} \\
Improvement ratio & \multicolumn{2}{c}{$2.21\times$} \\
\bottomrule
\end{tabular}
\end{table}

The hierarchical system produces patches on $53/300$ instances ($17.7\%$
patch generation rate) versus $24/300$ ($8.0\%$) for the baseline---a
$2.21\times$ improvement. Among the $49 + 20 = 69$ discordant pairs
(instances where the two conditions disagree), H wins $49$ ($71.0\%$) and B
wins $20$ ($29.0\%$). This asymmetry yields McNemar's test statistic
$\chi^2 = (|49 - 20| - 1)^2 / (49 + 20) = 11.36$ with $p = 6.0
\times 10^{-4}$, well below the $\alpha = 0.001$ threshold.

We additionally report 95\% confidence intervals on the patch rates using
the Wilson score interval \cite{wilson1927}:
\begin{align}
  \text{CI}_{95}(\rho_H) &= [13.4\%,\; 22.6\%], \\
  \text{CI}_{95}(\rho_B) &= [5.3\%,\; 11.7\%].
\end{align}
The two intervals are non-overlapping, confirming the significance of the
difference.

% --------------------------------------------------------
\subsection{Per-Repository Breakdown}
\label{sec:results:per_repo}
% --------------------------------------------------------

Table~\ref{tab:per_repo} reports patch generation rates and win/loss counts
broken down by repository. Three observations are noteworthy.

\begin{table}[t]
\centering
\caption{Per-repository patch generation results. H wins: instances where
only H patched. B wins: instances where only B patched.}
\label{tab:per_repo}
\small
\begin{tabular}{lrrrr}
\toprule
\textbf{Repository} & \textbf{N} & \textbf{H patches} & \textbf{B patches} & \textbf{H wins / B wins} \\
\midrule
sympy       & 77  & 0 (0.0\%)   & 0 (0.0\%)   & 0 / 0 \\
django      & 114 & 24 (21.1\%) & 17 (14.9\%) & 20 / 13 \\
scikit-learn & 23 & 14 (60.9\%) & 4 (17.4\%)  & 11 / 1 \\
sphinx-doc  & 16  & 10 (62.5\%) & 0 (0.0\%)   & 6 / 0  \\
pytest-dev  & 17  & 5 (29.4\%)  & 3 (17.6\%)  & 3 / 1  \\
matplotlib  & 23  & 4 (17.4\%)  & 2 (8.7\%)   & 2 / 0  \\
mwaskom     & 4   & 3 (75.0\%)  & 1 (25.0\%)  & 3 / 1  \\
psf         & 6   & 2 (33.3\%)  & 1 (16.7\%)  & 1 / 0  \\
pydata      & 5   & 3 (60.0\%)  & 0 (0.0\%)   & 2 / 0  \\
astropy     & 6   & 2 (33.3\%)  & 0 (0.0\%)   & 1 / 0  \\
pylint-dev  & 6   & 0 (0.0\%)   & 0 (0.0\%)   & 0 / 0  \\
pallets     & 3   & 0 (0.0\%)   & 0 (0.0\%)   & 0 / 0  \\
\midrule
\textbf{Total} & \textbf{300} & \textbf{53 (17.7\%)} & \textbf{24 (8.0\%)} & \textbf{49 / 20} \\
\bottomrule
\end{tabular}
\end{table}

\paragraph{Finding 1: sympy constitutes a complete failure for both conditions.}
All 77 sympy instances produce zero patches from either condition, with
sub-second elapsed time---indicating a systematic repository setup failure
(likely a shallow clone or checkout issue preventing \texttt{git diff} from
executing on the sympy commit history). Excluding sympy, the non-sympy
patch rates are $53/223 = 23.8\%$ (H) and $24/223 = 10.8\%$ (B), with
McNemar $\chi^2 = 11.36$ (unchanged, since zero discordant pairs from sympy
contribute to the statistic).

\paragraph{Finding 2: The largest improvement is on structurally complex codebases.}
Scikit-learn (H: 60.9\%, B: 17.4\%) and sphinx-doc (H: 62.5\%, B: 0.0\%)
show the most dramatic hierarchical advantage. These repositories have
modular, deeply nested package structures requiring multiple file reads
and cross-file reasoning to localize a bug---precisely the regime where
the monolithic baseline's context is most likely to be exhausted.

\paragraph{Finding 3: The baseline is not dominated.}
The baseline achieves 20 exclusive wins (instances where B patched but H
did not), demonstrating that the hierarchical system does not uniformly
dominate. We hypothesize that the overhead of Domain Manager planning
(consuming tokens on inter-level prompts and synthesis) occasionally
exhausts the per-instance timeout on simpler instances where the baseline
resolves the issue in a single direct \texttt{str\_replace} call.

% --------------------------------------------------------
\subsection{Efficiency and Resource Metrics}
\label{sec:results:efficiency}
% --------------------------------------------------------

Table~\ref{tab:efficiency} reports the aggregated time and token statistics.

\begin{table}[t]
\centering
\caption{Efficiency metrics across all 300 instances. Values are
mean $\pm$ standard deviation over instances with $>0$ tokens generated.}
\label{tab:efficiency}
\begin{tabular}{lcc}
\toprule
\textbf{Metric} & \textbf{Hierarchical} & \textbf{Baseline} \\
\midrule
Wall time (s)       & $94.8 \pm 73.9$ & $5.2 \pm 7.3$ \\
Total tokens        & $55{,}558 \pm \sigma_H$ & $7{,}177 \pm \sigma_B$ \\
Tokens / second     & $440.4$   & $1{,}615.0$ \\
Avg. tokens / agent & $3{,}388$ & $7{,}177$    \\
\midrule
Wall time | patched (H) & $146.7\,\text{s}$ & --- \\
Wall time | not patched (H) & $83.6\,\text{s}$ & --- \\
\bottomrule
\end{tabular}
\end{table}

Several observations follow from Table~\ref{tab:efficiency}:

\paragraph{Token speed paradox.} The hierarchical system generates tokens
at $440\,\text{tokens/s}$ versus the baseline's $1{,}615\,\text{tokens/s}$.
This apparent paradox---more expensive system, slower token rate---arises
because throughput is measured as total tokens / total wall time. In the
hierarchical system, the vLLM engine processes multiple sequential short
bursts (one per agent invocation) with inter-burst latency from tool
execution; the baseline makes fewer, longer generation calls with higher
sustained throughput. The true comparison is \emph{effectiveness per token},
not tokens per second.

\paragraph{Per-agent context efficiency.} Despite consuming 7.7$\times$
more total tokens per instance, the hierarchical system generates only
$3{,}388$ tokens per individual agent invocation, versus $7{,}177$ tokens
for the single baseline agent---a 52.8\% reduction in per-agent context
load. This directly validates the context isolation hypothesis: each agent
operates well within its effective context capacity.

\paragraph{Success correlates with token budget.} Hierarchical instances
that result in a patch consume significantly more tokens on average ($68{,}943$
tokens) than instances that fail to patch ($35{,}141$ tokens). This is
consistent with the hypothesis that successful instances require more
\emph{exploration} (more tool calls, more agent iterations) before
converging on the correct fix.
